\chapter{Background}\label{Background}
This chapter summarises the background required to contextualise the proposed coherent
receiver carrier-recovery architecture. It focuses on (i) short-reach IM/DD links versus
coherent detection, (ii) phase-related impairments (carrier frequency offset and laser
phase noise), and (iii) representative DSP-based carrier recovery methods used in coherent
receivers \cite{Savory2010JSTQE,Kuschnerov2009JLT}.

\section{Coherent optical communication systems}
Traditional short-reach optical interconnects in data centres have often favoured
intensity-modulation/direct-detection (IM/DD) transceivers because they enable simple and
power-efficient receiver front-ends \cite{Wartak2013OpticalReceivers}. In IM/DD systems,
information is conveyed by modulating the received optical intensity (or optical power).
At the receiver, a photodiode converts incident optical power into an electrical
photocurrent (square-law photodetection), so the receiver directly observes intensity but
does not directly access the optical carrier phase \cite{Wartak2013OpticalReceivers,KahnHo2004JSTQE}.

IM/DD links are attractive for short distances and moderate data rates; however, as symbol
rates increase, fibre chromatic dispersion becomes increasingly detrimental because group
delay varies with optical wavelength, causing waveform distortion and frequency-selective
fading for intensity-modulated signals. In addition, the achievable spectral efficiency
(information rate per unit bandwidth, in b/s/Hz) is fundamentally constrained when only
intensity is detected, because direct detection effectively exploits fewer degrees of
freedom than coherent detection \cite{KahnHo2004JSTQE}.
 
Coherent optical communication systems overcome these limitations by detecting both the
amplitude and phase of the optical field and enabling advanced modulation formats such as
QPSK and square M-QAM. Coherent detection also supports digital compensation of linear
channel impairments (e.g., chromatic dispersion and polarization effects) in the electrical
domain \cite{KahnHo2004JSTQE,Savory2010JSTQE,Kuschnerov2009JLT}. Recent work on data-centre
optics evaluates when coherent detection becomes favourable relative to IM/DD as per-lane
rates scale, highlighting the trade-off between coherent DSP power/complexity and improved
spectral efficiency and impairment tolerance \cite{Zhou2020JLT}.

A typical intradyne coherent receiver combines the received optical field with a local
oscillator (LO) laser in a 90$^\circ$ optical hybrid. The hybrid generates orthogonal
interference terms that are detected by balanced photodiode pairs, yielding electrical
in-phase (I) and quadrature (Q) components. Balanced detection improves sensitivity to the
desired beating term and suppresses common-mode intensity noise \cite{Savory2010JSTQE,Kikuchi2016JLT}.

After analog-to-digital conversion, the coherent receiver typically performs several DSP
operations, including chromatic dispersion compensation, timing recovery, adaptive
equalization, carrier frequency recovery and carrier phase recovery \cite{Savory2010JSTQE,Kuschnerov2009JLT}.
In many DSP chains, dispersion compensation and polarization demultiplexing/equalization
are performed prior to carrier phase estimation because commonly used blind equalizers can
converge without an explicit carrier-phase reference; in practice, frequency-offset
estimation may be placed either before or after equalization depending on acquisition range
and implementation constraints \cite{Savory2010JSTQE,Kuschnerov2009JLT}.

Historically, optical coherent receivers often relied on analogue phase-locked loops (PLLs)
for carrier synchronisation. In modern DSP-based coherent receivers, carrier recovery is
commonly implemented digitally using feed-forward and/or decision-directed algorithms,
which avoids tight analogue loop dynamics and integrates naturally with polarization-diverse
digital processing \cite{Savory2010JSTQE,IpKahn2007JLT,Pfau2009JLT}.

\section{Phase noise and frequency offset}
The main phase-related impairments in coherent optical communication systems are carrier
frequency offset (CFO) and phase noise. Both arise because the transmitter and local
oscillator lasers are independent devices and therefore do not share an identical optical
frequency or phase reference \cite{Kikuchi2016JLT}.

Carrier frequency offset occurs because the instantaneous optical frequency of the
transmitter laser differs from that of the LO laser. After coherent detection and
digitisation, CFO appears as an effective baseband frequency offset and produces a
deterministic phase ramp across samples \cite{Kikuchi2016JLT,Leven2007LPT}. Mathematically,
this phase rotation can be expressed as
\begin{equation}
y[k] = s[k] e^{j(\theta[k] + 2\pi \Delta f T_s k)} + \eta[k]
\end{equation}
where $s[k]$ is the transmitted symbol, $\theta[k]$ is the phase noise, $\Delta f$ is the
carrier frequency offset, $T_s$ is the symbol period and $\eta[k]$ is additive noise
\cite{Kikuchi2016JLT,RoudasCoherentSystems}.

The term $2\pi \Delta f T_s k$ corresponds to a linear phase ramp. A larger frequency offset
or a larger symbol duration leads to faster constellation rotation. This is particularly
relevant for block-based carrier recovery systems, because the total accumulated phase
change within one block influences the probability of phase ambiguities and cycle slips
\cite{Pfau2009JLT}.

Besides frequency offset, coherent optical systems are also affected by phase noise. Laser
phase noise is commonly modelled as a discrete-time Wiener process (random walk) \cite{RoudasCoherentSystems}:
\begin{equation}
\theta[k] = \theta[k-1] + \Delta \theta[k]
\end{equation}
where $\Delta\theta[k]$ is a zero-mean Gaussian increment whose variance is proportional to
the combined (beat) linewidth and the sampling interval \cite{RoudasCoherentSystems,Pfau2009JLT}.

For the systems considered in this thesis, the residual CFO after coarse frequency recovery
is especially important. While a separate frequency recovery stage can compensate most of
the CFO, a small leftover frequency error remains due to the finite resolution of practical
estimators \cite{Leven2007LPT,Selmi2009ECOC}. This residual offset causes a nearly linear
phase drift across each DSP block.
\missingfigure{figure illustrating the near-lineair approximation
of the total noise process}

The phase recovery system must therefore not only compensate random phase noise caused by
laser linewidth, but also account for remaining within-block phase slope. This motivates
the channel model used throughout this thesis, in which the received symbols are impaired by:
\begin{itemize}
    \item additive white Gaussian noise,
    \item a residual carrier frequency offset,
    \item Wiener phase noise caused by laser linewidth.
\end{itemize}

A coarse FFT-based frequency recovery algorithm is often used before phase recovery.
FFT-based CFO estimation provides wide acquisition range at modest complexity
\cite{Leven2007LPT,Selmi2009ECOC}. For M-PSK constellations, an M-th power nonlinearity can
suppress the data modulation and expose a strong carrier-related spectral line that enables
peak-based CFO estimation; related techniques are also adapted for square QAM via
partitioning, pilots, or data-aided processing \cite{Rasmussen2011SPIE}.

The estimated frequency offset can then be compensated digitally. However, the accuracy of
this estimate depends on the FFT size. A finite FFT resolution leads to a remaining
frequency estimation error. Since the FFT bin spacing equals $F_s/N_{\mathrm{FFT}}$, a
simple bin-centred peak estimate implies an error on the order of half a frequency bin:
\begin{equation}
\Delta f_e \leq \frac{F_s}{2N_{\mathrm{FFT}}}
\end{equation}
\cite{AnalogDevicesDFT,AgilentNoiseGuide}.

\missingfigure{add figure explaining frequency error due to finite FFT leads to
offloading to phase recovery stage}
This remaining frequency error is then passed on to the phase recovery stage. If the
residual offset is too large, the phase recovery block may no longer be able to track the
constellation rotation correctly, which increases the probability of cycle slips
\cite{Pfau2009JLT,Kuschnerov2009JLT}.

\section{Carrier recovery in DSP chains}
Carrier recovery is the DSP stage responsible for compensating the phase rotations
introduced by residual CFO and phase noise. Since these impairments have different
characteristics, coherent receivers commonly separate carrier recovery into (i) frequency
offset estimation/compensation and (ii) carrier phase estimation \cite{Savory2010JSTQE,Kuschnerov2009JLT}.

Frequency recovery estimates a single CFO parameter over an observation interval and
therefore benefits from aggregating many samples (or blocks) to improve estimator variance;
it is not typically meaningful to produce an independent CFO estimate per symbol. By
contrast, carrier phase recovery must track a time-varying phase process (Wiener phase
noise plus residual CFO), motivating symbol-rate or block-rate phase estimation depending
on implementation parallelism \cite{Leven2007LPT,Pfau2009JLT}.

A major challenge is that modern coherent DSP systems process many symbols in parallel in
order to reduce clock speed requirements. As a result, carrier phase is often estimated
only once per processing block instead of once per symbol \cite{Kuschnerov2009JLT}. Although
this significantly reduces computational complexity, it also increases sensitivity to phase
drift within the block. If the total phase change over a block becomes too large, the phase
recovery algorithm may unwrap to an incorrect $\pi/2$-shifted hypothesis (cycle slip) for
square QAM constellations \cite{Pfau2009JLT}.

If a block contains $B$ symbols and the residual frequency offset is $\Delta f_e$, the total
phase drift across the block can be approximated by
\begin{equation}
\Delta \phi_{\mathrm{block}} \approx 2 \pi \Delta f_e B T_s
\end{equation}
To reduce the probability of cycle slips, a practical design rule is to keep the intra-block
phase drift well below the midpoint between neighbouring $\pi/2$-ambiguous hypotheses, often
expressed as $|\Delta \phi_{\mathrm{block}}| \lesssim \pi/4$ \cite{Pfau2009JLT}.

\section{Blind Phase Search}
Blind Phase Search (BPS) is a widely used feed-forward carrier phase estimation method for
coherent square M-QAM. Its advantages include pilot-free operation, modulation-format
compatibility, and strong tolerance to laser phase noise when sufficient test phases and
averaging are used \cite{Pfau2009JLT,Yang2018OE}.

The principle of BPS is straightforward. A block of received symbols is rotated over a set
of candidate test angles. For each candidate angle, the rotated symbols are sliced to the
nearest ideal constellation points and an error metric is computed. The candidate angle
that minimises the total error is selected as the estimated carrier phase \cite{Pfau2009JLT}.

\missingfigure{block diagram of working bps principle + testangles on constellationdiagram}

For a received block $z[k]$, the BPS metric for a test phase $\phi$ can be written as
\begin{equation}
e(\phi) = \sum_{k=0}^{B-1} \left| D\left(z[k] e^{-j\phi}\right) - z[k] e^{-j\phi} \right|^2
\end{equation}
where $D(\cdot)$ denotes the decision operation that maps a rotated symbol to its nearest
constellation point. The phase estimate is then given by the test angle that minimises
$e(\phi)$.

For square QAM constellations, the objective function in BPS is periodic with period $\pi/2$
due to constellation rotational symmetry. Therefore, it suffices to search any interval of
length $\pi/2$; a common centred choice is $[-\pi/4,\pi/4]$ \cite{Pfau2009JLT}. However, a
large number of test phases is usually required to obtain sufficient estimation accuracy,
making BPS computationally expensive in high-throughput FPGA implementations
\cite{Pfau2009JLT,Yang2018OE}.

Another limitation of blockwise BPS is that it estimates one average phase per block. When
residual CFO induces significant intra-block phase drift, a single blockwise estimate can be
insufficient and may increase cycle-slip probability \cite{Pfau2009JLT}. This thesis therefore
extends the traditional blockwise BPS approach by also estimating the phase spread across the
block, as explained in Chapter~\ref{Proposed phase recovery architecture}.

\section{CORDIC}
The Coordinate Rotation Digital Computer (CORDIC) algorithm implements vector rotations and
elementary functions using iterative shift-add operations. Its attractiveness in FPGA/ASIC
implementations stems from replacing general multipliers with pipelinable add/subtract and
bit-shift datapaths \cite{Volder1959TEC,Meher2009TCASI}.

In the rotation mode of CORDIC, a vector is iteratively rotated by a sequence of predefined
micro-rotations. The iterative equations can be written as \cite{Volder1959TEC,Meher2009TCASI}:
\begin{equation}
x_{i+1} = x_i - d_i y_i 2^{-i}
\end{equation}
\begin{equation}
y_{i+1} = y_i + d_i x_i 2^{-i}
\end{equation}
\begin{equation}
z_{i+1} = z_i - d_i \arctan(2^{-i})
\end{equation}
where $d_i \in \{-1, +1\}$ determines the direction of rotation at iteration $i$.
After a sufficient number of iterations, the vector is rotated by the desired angle with
good precision \cite{Meher2009TCASI}.

In this thesis, CORDIC is used to rotate received symbols by the estimated carrier phase.
The implementation leverages an FPGA IP core; design effort therefore focuses on selecting
fixed-point formats, pipeline depth/latency, and interface timing to meet throughput and
numerical-accuracy requirements \cite{AMD2021PG105}.
  